% Part: proof-theory
% Chapter: normalization
% Section: introduction

\documentclass[../../../include/open-logic-section]{subfiles}

\begin{document}

\olfileid{pt}{nor}{seg}

\olsection{Segments and Cuts}

!!^a{introduction} rule followed by !!a{elimination} rule is a detour
in !!a{proof}, and to normalize !!a{proof} is to eliminate such
detours. However, the \Elim{\lor} and \Elim{\lexists} make even the
definition of ``detour we want to eliminate'' a bit complicated. The
problem is this. Because of these rules, it's not guaranteed that in a
detour the !!{elimination} rule \emph{immediately} follows the
!!{introduction} rule. There could be an \Elim{\lor} or
\Elim{\lexists} in between, e.g.,
\[
\AxiomC{}
\DeduceC{$\lexists[x][!C(x)]$}
\AxiomC{$\Discharge{!A}{x}$}
\DeduceC{$!B$}
\DischargeRule{\Intro{\lif}}{x}
\UnaryInfC{$!A \lif !B$}
\DischargeRule{\Elim{\lexists}}{y}
\BinaryInfC{$!A \lif !B$}
\AxiomC{}
\DeduceC{$!A$}
\RightLabel{\Elim{\lif}}
\BinaryInfC{$!B$}
\DisplayProof
\]
In fact, any number of \Elim{\lor} or
\Elim{\lexists} can separate the \Intro{\lif} from the \Elim{\lif}, e.g.,
\[
\AxiomC{}
\DeduceC{$\lexists[x][!C(x)]$}
\AxiomC{}
\DeduceC{$!D \lor !E$}
\AxiomC{$\Discharge{!A}{x}$}
\DeduceC{$!B$}
\DischargeRule{\Intro{\lif}}{x}
\UnaryInfC{$!A \lif !B$}
\AxiomC{$\Discharge{!A}{x}$}
\DeduceC{$!B$}
\DischargeRule{\Intro{\lif}}{x}
\UnaryInfC{$!A \lif !B$}
\RightLabel{\Elim{\lor}}
\TrinaryInfC{$!A \lif !B$}
\DischargeRule{\Elim{\lexists}}{y}
\BinaryInfC{$!A \lif !B$}
\AxiomC{}
\DeduceC{$!A$}
\RightLabel{\Elim{\lif}}
\BinaryInfC{$!B$}
\DisplayProof
\]
This example shows that the same \Elim{\lif} might be involved in
multiple detours, even, as it follows more than one~\Intro{\lif}. To
account for all these possibilities, we define detours using certain
kinds of sequences of !!{formula} occurrences in an \Log{N1i}-proof
called \emph{segments}. In our example, it's a sequence of occurrences
of~$!A \lif !B$, where a minor premise of $\Elim{\lor}$ or
$\Elim{\lexists}$ follows its conclusion. In the example, there are
two such segments, each containing three occurrences of~$!A \lif !B$.
Here is the official definition.

\begin{defn}
A \emph{segment} in a \Log{N1i} !!{proof}~$\delta$ is a sequence of
occurrences $!A_1$, \dots, $!A_n$ of the same !!{formula}~$!A$ in $\delta$ where
\begin{enumerate}
  \item $!A_1$ is not the conclusion of \Elim{\lor} or \Elim{\lexists};
  \item $!A_n$ is not a minor premise of \Elim{\lor} or \Elim{\lexists};
  \item For $i = 1$, \dots,~$n-1$, $!A_i$ is a minor premise of
  \Elim{\lor} or \Elim{\lexists} and $A_{i+1}$ is the conclusion of
  that inference.
\end{enumerate}
The \emph{length} of that segment is~$n$.
\end{defn}

Not ever segment is a target for elimination (a detour). The ones that
are are called \emph{cuts}.

\begin{defn}
A \emph{cut} is a segment in which $!A_n$ is the major premise of
!!a{elimination} inference, and either $n=1$ and $!A_1 = !A_n$ is the
conclusion of !!a{introduction} inference or of~\FalseInt, or~$n>1$.
\end{defn}

Note that a cut segment need not begin with the conclusion of
!!a{introduction} rule. All that's required is that it ends with the
major premise of !!a{elimination} rule. However, a segment of
length~$1$ does have to be the conclusions of !!{introduction} rules
to count as a cut.

\begin{defn}
The \emph{rank} of a cut is~$\depth{!A}$. The \emph{cut
rank}~$\cutrank{\delta}$ of !!a{proof}~$\delta$ is the maximal rank of
cuts in~$\delta$. A cut is \emph{maximal} in~$\delta$ if its rank
is~$\cutrank{\delta}$, i.e., it has maximal rank. The \emph{cut length}
of~$\delta$ is the sum of the lengths of all its maximal cuts.
\end{defn}

\end{document}
